\section*{Chapter \ref{ch:s1:the-asset-managers-strategies} --- The Asset Managers' Strategies}

\begin{solution}{exo:s1:the-asset-managers-strategies:1}
By $1/\sqrt{2} = 0.707$. The table shows $4.84/2.96 = 1.64$ from 25 to 50 names and $1.64/0.81 = 2.02$ from 100 to 200: faster than the square root,
because the names added are ever smaller in the index and the weight left out shrinks faster than the number of names grows.
\end{solution}

\begin{solution}{exo:s1:the-asset-managers-strategies:2}
$0.064 \times 10 = 0.64$ basis points a year.
\end{solution}

\begin{solution}{exo:s1:the-asset-managers-strategies:3}
Tracking error measures how far the weights are from the index, which grows with the tilt; the return depends on what the tilted factor earns, which the
product does not control.
\end{solution}

\begin{solution}{exo:s1:the-asset-managers-strategies:4}
$3.46\%/\sqrt{252} = 0.218\%$, 21.8 basis points a day. Over a one-day transition the whole trade list is exposed for one day to the difference between the
old and new portfolios, whose daily standard deviation this is.
\end{solution}

\begin{solution}{exo:s1:the-asset-managers-strategies:5}
Between rebalances the held weights drift with returns while the optimal weights would change, listings enter and leave the index, and the ex-ante figure
is an average over rebalance dates of forecasts made with that day's exposures.
\end{solution}

\begin{solution}{exo:s1:the-asset-managers-strategies:6}
$0.03 \times 0.52 = 1.56\%$ a year. It earned 0.07\% because book-to-price measured with today's price loads negatively on the planted momentum, whose
premium offsets the value premium (chapter~6).
\end{solution}

\begin{solution}{exo:s1:the-asset-managers-strategies:7}
Cost plus two standard deviations is 71.2, 70.0, 72.8, 80.3 and 98.3 basis points for 1, 2, 3, 5 and 10 days: the planner chooses two days. More
aversion to risk shortens the transition.
\end{solution}

\begin{solution}{exo:s1:the-asset-managers-strategies:8}
A sample of the 50 largest names is a bet on large caps against the rest; with a tracking error near 3\%, ten years of 2\% a year is within about two
standard errors of zero, and the synthetic 50-name sample did the same with no skill. The purpose of sampling is low tracking error at low cost, not
excess return.
\end{solution}

\begin{solution}{pb:s1:the-asset-managers-strategies:1}
\begin{enumerate}
\item Holding every constituent at index weight; holding a subset weighted to match the index's risk.
\item Tracking error is unavoidable because of frictions, and \omterm{def:s1:index-rebalancing:fund}{index funds} on average beat active funds after expenses.
\item Cap weights drift with prices by themselves; only entries and exits need trades (6.4\% a year).
\item The market's true exposures (beta, industries, point-in-time styles) and specific risks; a QP over the largest $n$ names minimising ex-ante
tracking error, long-only.
\item 4.84, 2.96, 1.64, 0.81 and 0.32\% realised; 4.70, 2.81, 1.52, 0.71 and 0.26\% ex ante, for 25 to 400 names.
\item Their weights must be re-optimised as the index moves, and each re-optimisation moves more weight.
\item Nothing: a bet on large caps within its tracking error.
\item Passive funds benefit from less rigid rebalancing; enhanced funds rebalance earlier and more patiently and pay less.
\item A rules-based tilt toward a style; exposures 0.26, 0.52 and 1.03, tracking errors 1.57, 3.08 and 6.00\%, returns 0.06--0.07\%.
\item Value measured with today's price is short momentum in this market.
\item Take tracking error $\omega = \text{IR}/(2\lambda)$ for information ratio IR and aversion $\lambda$ to active variance.
\item Moving holdings between portfolios at the least shortfall, balancing impact against the risk of the unexecuted part.
\item \textbf{Named result.} \omterm{def:s1:the-asset-managers-strategies:replication}{Sampled replication} tracks the synthetic index to 4.84\%, 2.96\%, 1.64\%, 0.81\% and 0.32\% a year with 25, 50, 100, 200 and 400
names; the transition of 59\% of a \$2 billion fund costs 27.6 basis points in one day ($\pm$21.8) and 18.5 in three ($\pm$27.2), with cost plus one
standard deviation least at two to three days (45.6).
\item Impact limits speed; active risk limits patience.
\item Names common to both portfolios, or traded against other flows, change owner without market impact.
\item Larger ex-ante errors and a larger gap to realised tracking error.
\item The index as it was, announcement dates, dividends, cash flows and costs.
\item Sampled index replication.
\item Chapter~10 traded the index changes that \omterm{def:s1:index-rebalancing:fund}{index funds} must follow; here the funds' side.
\item Keeping a portfolio close to a target cheaply, and moving it without losing money on the way.
\end{enumerate}
\end{solution}

\begin{solution}{iq:s1:the-asset-managers-strategies:1}
Costs and fees, cash drag from flows and dividends, trading at prices other than the index's closing prices, sampling instead of \omterm{def:s1:the-asset-managers-strategies:replication}{full replication}, and
timing of trades around index changes.
\end{solution}

\begin{solution}{iq:s1:the-asset-managers-strategies:2}
Hold the largest names fully, then choose among the rest by an optimiser that matches industry, country and style exposures at the least tracking error,
with liquidity limits; or stratify by industry and size and pick representative names in each cell.
\end{solution}

\begin{solution}{iq:s1:the-asset-managers-strategies:3}
Build the trade list, cross common names and internal flows first, hedge interim exposure with futures if the transition is long, schedule the rest to
balance impact against active risk, and report implementation shortfall against the decision prices.
\end{solution}

\begin{solution}{iq:s1:the-asset-managers-strategies:4}
The provider's announcement with adds, deletes and weight changes, corporate actions, the effective date and time, expected trade sizes against
liquidity, and checks that the resulting holdings match the new index within tolerance.
\end{solution}

\begin{solution}{iq:s1:the-asset-managers-strategies:5}
Its exposure to the target factor and to others it picks up (a value product's momentum exposure), tracking error against the parent index, crowding in
the factor, and liquidity; worry when unintended exposures grow or the factor's correlation with others changes.
\end{solution}

\begin{solution}{iq:s1:the-asset-managers-strategies:6}
$\sum_{d=0}^{N-1}(1 - d/N)^2 \approx N\int_0^1(1 - x)^2dx = N/3$, so the risk is about $\sigma_a\sqrt{N/3}$. Minimising $c_0N^{-1/2} + \lambda\sigma_a
\sqrt{N/3}$ gives $-\tfrac12 c_0 N^{-3/2} + \tfrac{\lambda\sigma_a}{2\sqrt3}N^{-1/2} = 0$, so $N = \sqrt3\,c_0/(\lambda\sigma_a)$.
\end{solution}
